% !TEX root = ../main.tex
\section{Assumptions and Methodology}
\label{sec:methodology}

\plantodo{S1}{State the working assumptions carried through the project --
see README.md "Working Assumptions" and
\texttt{projects/nas-sos-capstone/index.md} "Working assumptions (guardrails)"
(e.g., emphasis on structure/behavior/interfaces/traceability rather than
optimization math; decision-support/optimization treated as a capability
inside the architecture, not the whole subject).}

\plantodo{S6}{Describe the methodology for exploring and comparing
alternative system decompositions (organization/authority-based,
lifecycle/process-based, physical-system, information-flow,
decision-authority, intent-to-trajectory) and the criteria used to evaluate
them (what each exposes/obscures, where boundaries are ambiguous, which
best supports optimization-impact analysis).}

\plantodo{S10}{Describe the MBSE methodology used to build the architecture
in SysML v2 (textual notation authored in Syside Modeler in VS Code, with
diagrams as generated views): package organization, view types produced (BDD,
IBD, activity, sequence, requirements), and the traceability approach
(stakeholder needs -> objectives -> requirements -> systems -> activities ->
information exchanges -> aircraft behavior). Note per
\texttt{workflows/update-architecture.md} (decision D-006) that the SysML v2
text is the source of truth; Cameo is not used as the source of truth.}

\plantodo{S11}{Describe the optimization-study methodology: scenario
selection criteria, system-state and decision-variable definitions,
constraint definitions, the multi-objective formulation
$J = w_1 J_1 + w_2 J_2 + \dots + w_n J_n$, normalization approach for
objectives with different units/scales, and the stakeholder weighting
scheme. This section should be detailed enough that the scenario and
formulation could be reproduced.}
